\documentclass[../../main.tex]{subfiles}
\begin{document}

\begin{center}
  {\fontsize{48pt}{58pt}\selectfont \textbf{1964}\normalsize\textbf{年度}}
\end{center}

\begin{table}[htb]
  \centering
  \caption{出題テーマと難易度}
  \begin{tabular}{|c|p{9cm}|c|}
    \hline
    \textbf{問題番号} & \textbf{テーマ} & \textbf{難易度} \\
    \hline
    第1問 & 三次方程式の解と係数の関係 & \\
    \hline
    第2問 & 放物線と内分点の軌跡 & \\
    \hline
    第3問 & 四角錐の切断面積 & \\
    \hline
    第4問 & 四辺形の頂点を繰り返し中点にとる図形 & \\
    \hline
    第5問 & 双曲線上の2点間の距離の最小化 & \\
    \hline
    第6問 & 三次関数の非負条件と定積分の最大・最小 & \\
    \hline
  \end{tabular}
\end{table}

\begin{table}[htb]
  \centering
  \caption{解答}
  \renewcommand{\arraystretch}{1.4}
  \begin{tabular}{|c|c|p{9cm}|}
    \hline
    \textbf{問題番号} & \textbf{小問} & \textbf{答え} \\
    \hline
    第1問 & -- & $a^3+b^3+c^3=z^3+3yz+3x$ \\
    \hline
    第2問 & -- & $P=(-6,3)$，$m:n=2:1$ \\
    \hline
    第3問 & -- & $350\,\mathrm{cm^2}$ \\
    \hline
    第4問 & -- & $P_1=(\dfrac{2}{3},\dfrac{8}{5})$，$P_2=(\dfrac{1}{3},\dfrac{4}{5})$，$P_3=(\dfrac{2}{3},\dfrac{2}{5})$，$P_4=(\dfrac{4}{3},\dfrac{6}{5})$．面積は$\dfrac{3}{5}$ \\
    \hline
    第5問 & (1) & $(a^{2/3}+a^{-2/3})^{3/2}$ \\
    \cline{2-3}
    & (2) & $a=3^{3/4}$ \\
    \hline
    第6問 & -- & 最大のとき$(a,b)=(-6,9)$，最小のとき$(a,b)=(3,0)$ \\
    \hline
  \end{tabular}
\end{table}

\paragraph{試験形態}
6問．試験時間は要確認．

\end{document}
